\section{Problem Formulation}
\label{sec:formulation}

\paragraph{Skills, positions and representations.}
Let $x$ be a question, $s$ its correct skill, and $\tilde{s}$ an unrelated skill
with the same token length. For a fixed autoregressive model of hidden width
$D$, let $h^L_{\mathrm{skill}}(p)$, $h^L_{\mathrm{wrong}}(p)$ and
$h^L_{\mathrm{no}}(p)$ denote residual-stream hidden states at layer $L$ and prompt position
$p$ with $s$, with $\tilde{s}$ and without a skill. At the final prompt position,
we decompose the total displacement as
\begin{equation}
\underbrace{h_{\mathrm{skill}}-h_{\mathrm{no}}}_{\tvec\ \text{(total)}}
=
\underbrace{h_{\mathrm{wrong}}-h_{\mathrm{no}}}_{\gvec\ \text{(presence)}}
+
\underbrace{h_{\mathrm{skill}}-h_{\mathrm{wrong}}}_{\dvec\ \text{(content)}}.
\label{eq:decomp}
\end{equation}
We omit layer and position indices in Eq.~\ref{eq:decomp} for readability.
``Presence'' denotes the change induced by the chosen wrong skill; ``content''
denotes the additional change induced by replacing it with the correct skill.
This algebraic decomposition does not assume independent underlying factors.
We test each component by adding it to the corresponding state in a no-skill run.

Across the skill's $m$ token positions $S$, the content-specific difference
is a matrix $\dvec_S^L=H^L_{\mathrm{skill}}(S)-H^L_{\mathrm{wrong}}(S)
\in\mathbb{R}^{m\times D}$. We use a run containing the wrong skill as the
\emph{receiver}, and the correct-skill run as the \emph{donor}. Injection at dose
$\alpha$ sets
\begin{equation}
H^L_{\mathrm{patched}}(S)=H^L_{\mathrm{wrong}}(S)+\alpha\dvec_S^L.
\label{eq:inject}
\end{equation}
At $\alpha=1$, the donor's states replace the receiver's across the span;
other states remain unchanged at that layer. Computation then continues.
Matching lengths preserves positions and their encodings. Span interventions
change many more states than final-position interventions, which limits direct
comparisons.

\paragraph{Behavioural recovery and evaluation subsets.}
We score generated answers and measure the fraction of the donor's accuracy
gain recovered by an intervention:
\begin{equation}
\rho(\mathrm{arm})=
\frac{\mathrm{acc}(\mathrm{arm})-\mathrm{acc}(\mathrm{receiver})}
     {\mathrm{acc}(\mathrm{donor})-\mathrm{acc}(\mathrm{receiver})}.
\label{eq:rho}
\end{equation}
Here, an arm is an intervention condition, and the donor baseline uses the
correct skill in context. Thus $\rho=1$ recovers the donor's observed gain,
and values above one indicate greater accuracy than the donor. We distinguish rescued
items (R: wrong without the skill, correct with it), persistent failures (F:
wrong under both), kept items (K: correct under both), and broken items (B:
correct without, wrong with). Main causal comparisons use rescued items. For
span interventions, we further retain only skill groups for which the
wrong-skill receiver solves none of the sampled items. We apply this screen
separately within each run, model and task, without using intervention outcomes.
The estimates are conditional on baseline failure in the sampled group;
they do not measure recovery over the full benchmark.

\paragraph{Rank and depth as intervention variables.}
To test how many directions support recovery, we subtract the mean across
positions, truncate the resulting matrix, and restore the mean:
\begin{equation}
\dvec^{(k)}=\mathbf{1}\bar{\dvec}+\sum_{i=1}^{k}\sigma_i u_i v_i^\top,
\qquad
\dvec-\mathbf{1}\bar{\dvec}=\sum_i\sigma_i u_i v_i^\top.
\label{eq:rank}
\end{equation}
Here, $\bar{\dvec}\in\mathbb{R}^{1\times D}$ is the mean state difference,
repeated at every position. The singular values $\sigma_i$ are ordered from
largest to smallest, with $u_i$ and $v_i$ their associated directions. Thus $k$
counts retained directions after centring, and the complete injected matrix has
rank at most $k+1$. A bottom-$k$ control keeps the $k$ smallest singular
directions and the same mean. We also vary the intervention layer. Span
replacement tests whether states at that layer suffice to transfer the effect;
blocking attention to the span from that layer onward tests whether subsequent
computation still depends on reading it.

